\documentclass[Helvetica]{reporte}
\usepackage[utf8]{inputenc}
%\usepackage{natbib}
%\usepackage{apacite}

%%%%%%%%%%%%%%%%%%%%%%%%%% Datos de la materia %%%%%%%%%%%%%%%%%%%%%%%%%%
\Materia {Sistemas Fotónicos}
\Seccion {D03}
\Profesor {Guillermo García Torales}
%%%%%%%%%%%%%%%%%%%%%%%%%% Datos del alumno %%%%%%%%%%%%%%%%%%%%%%%%%%
\Apellidos {XXXX XXX}
\Nombre{XXX XXXX}
\Codigo {2XXXXXXX24}
%%%%%%%%%%%%%%%%%%%%%%%%%% Datos de la práctica %%%%%%%%%%%%%%%%%%%%%%%%%%
\Numero {1}
\Titulo {Modulación con Puente Wheatstone}

%%%%%%%%%%%%%%%%%%%%%%%%%% inicio del documento%%%%%%%%%%%%%%%%%%%%%%%%
\iniciadocumento
 
\section{Introducción}
La introducción debe contener una explicación breve que relacione la problemática a resolver, el objetivo y los antecedentes del trabajo a realizar.
\section{Marco teórico}
Aquí se deberá añadir los fundamentos teóricos en los cuales se sustenta la práctica. debiendo colocar las referencias y ecuaciones apropiadas. 

Por ejemplo: 
\cite{rashid} \cite{boylestad} \cite{rashidelectronica}\cite{boylestadE} nueva cita \cite{chapman}
\newline
Ejemplo ecuación:
ecuación en linea $a=b+c$ \\%doble diagonal invertida significa cambio de linea
ecuación en una linea aparte $$a=a+b$$
Ecuación numerada
\begin{equation}
L=\lim_{t\to c}\sin (\omega t)
\label{eq:limite}
\end{equation}
\begin{equation}
a=b+c
\label{eq:a}
\end{equation}
A la ecuación numerada se le puede llamar usando \eqref{eq:a}

\section{Metodología y materiales}
 En esta sección se debe describir el procedimiento de cada una de las etapas relevantes de la práctica. Explicando el funcionamiento mediante una descripción detallada y basada en Figuras, calculos, tablas, gráficas, fotografias, simulaciones, etc.
 
 \begin{enumerate}
 \item Etapa 1
 \end{enumerate}
 
\begin{itemize}
\item Cálculos
\item Diagramas
\item Fotografías 
\item Tablas (Listas de materiales y equipo utilizado) \end{itemize}

\begin{enumerate}
\item Etapa 2
\end{enumerate}

\begin{itemize}
\item Cálculos
\item Diagramas
\item Fotografías 
\item Tablas (Listas de materiales y equipo utilizado) \end{itemize}

%%%%%%%%%%%%%%%%%%%%%%%%%%%
FIGURAS
%%%%%%%%%%%%%%%%%%%%%%%%%%
\begin{figure}[H]  %iniciar el entorno figura
	\centering	%Centra la figura
	\includegraphics[width=16cm]{Diagramas/diagrama1.png}
	\caption{Descripción de la figura.}
	\label{fig:diagrama1}
\end{figure}
Todas las figuras deben ser descritas en el texto: por ejemplo, la figura \ref{fig:diagrama1} muestra

\Tabla{Lista de materiales}%Titulo de la tabla
{c|l} % describir las columnas l-a la izquierda r-a la derecha y c- centrado, |- linea intermedia entre columnas
{Cantidad &Descripción\\% & separa columna \\ - cambio de linea
\hline %linea horizontal en la tabla
2& Resistencias de 220$\Omega$\\
4& Diodos 1N4007 (ver \ref{hoja:1N4007})\\ % en elementos electrónicos es necesario referenciar las hojas de datos
1 & Transistor 2N2222 (ver \ref{hoja:2N2222})\\
}
{tab:materiales} %etiqueta
Tambien las tablas deben referenciarse en el texto \ref{tab:materiales}


\lstinputlisting[language=C++,style=c,caption={Ejemplo de código fuente en C}]{Codigos/hello.c}

\lstinputlisting[language=Arduino,caption={Ejemplo de código fuente de arduino}]{Codigos/codigo.ino}



\section{Resultados}
 En los resultados se deberán incluir tablas y gráficas donde se registren las mediciones realizadas, así como su análisis en orden cero. Además se pueden mostrar fotografías del prototio funcionando. Se debe incluir una discusión sobre los resultados y explicar que es lo relevante de la técnica o medición realizada.
 
 
\begin{figure}[H]
\centering
\begin{tikzpicture}%entorno para graficas
\pgfplotsset{domain=0:180} %eje x
\begin{axis}[ cycle list name=color list] %inicia la grafica
\addplot [color=blue, mark=none, dashed]{2*5*(1.5151-1)*cos(x/2)};%[opciones]{formmula}
\addplot table [color=blue, mark=none, x=omega, y=a5, col sep=comma ] {Datos/Risleyparaxialvsexacta.csv};%[x=columnadel eje de las x, y= columna del eje de las y]{ubicacion de la tabla/tabla completa}

\addplot [blue, mark=none, dashed]{2*10*(1.5151-1)*cos(x/2)};
\addplot table [ mark=none, x=omega, y=a10, col sep=comma ] {Datos/Risleyparaxialvsexacta.csv};

\addplot [blue, mark=none, dashed]{2*15*(1.5151-1)*cos(x/2)};
\addplot table [ mark=none, x=omega, y=a15, col sep=comma ] {Datos/Risleyparaxialvsexacta.csv};

\addplot [blue, mark=none, dashed]{2*20*(1.5151-1)*cos(x/2)};
\addplot table [mark=none, x=omega, y=a20, col sep=comma ] {Datos/Risleyparaxialvsexacta.csv};
%\addlegendentry{Desviación exacta}
%\addlegendentry{Aproximación paraxial}
\end{axis}
\end{tikzpicture}
\caption{Descripción de la figura.  }
\label{fig:paraxial}
\end{figure}

\section{Conclusiones}
En esta sección no se agrega nada nuevo, solo se vuelve a presentar de forma sucinta lo más sobresaliente. No escribir agradecimientos ni notas personales.



%%%%%%%%%%%%%%%%%%%%%%%%%%%%%%%% Apéndices %%%%%%%%%%%%%%%%%%%%%%%%%%%%%%%%%%%%%%%%%%%%%%
\newpage
\appendix	%Aquí inician los apéndices
\section{Hojas de datos}
\secttoc %agrega una lista de las hojas de datos

\incluirpdf{Hojas_de_datos/1N4007.pdf}{1}{2}{Diodo rectificador 1N4007}{hoja:1N4007}%{dirección del pdf}{pagina inicial}{pagina final/se pueden especificar mas intervalos de paginas separados por comas}{Titulo}{etiqueta}

\incluirpdf{Hojas_de_datos/2N2222.pdf}{1}{3}{Transistor de switcheo 2N2222}{hoja:2N2222}

%\bibliographystyle{apa} 
\bibliographystyle{IEEEtran}	
%\bibliographystyle{apacite}
\bibliography{Biblio}
\end{document}
